\phantomsection\addcontentsline{toc}{section}{Biomimetics}
\ECchapter{19}{Biomimetics}{How can engineering learn from living systems?}{ECPurple}

\ECObjectives{
\begin{itemize}
\item Distinguish inspiration from direct imitation of a biological model.
\item Explain how structure, surface and control principles can be transferred to engineering.
\item Evaluate a biomimetic solution using measurable performance and life-cycle criteria.
\end{itemize}}

\begin{multicols}{2}
\begin{engineeringnews}
Biomimetic research is influencing adhesives, lightweight structures, ventilation systems and robot
locomotion. Nature provides useful strategies, but engineering success requires abstraction,
modelling and testing rather than simply copying the appearance of an organism.
\end{engineeringnews}

\begin{ecarticle}
Biomimetics examines biological functions and transfers useful principles to engineered systems. The
process usually begins with a technical problem such as reducing drag, gripping a surface or
building a lightweight structure. Engineers then search for organisms that solve a comparable
functional problem under similar constraints.

The biological solution must be understood at the relevant scale. A bird wing is not only a curved
surface; its performance depends on flexible feathers, muscles, sensing and active control. Directly
copying its shape may therefore fail. Engineers abstract the underlying principle, create a model and
adapt it to available materials and manufacturing methods.

Surface engineering provides many examples. Gecko feet adhere through millions of microscopic hairs
that create weak intermolecular forces over a large contact area. Artificial dry adhesives reproduce
the hierarchical structure without using glue. Shark skin has rib-like textures that influence the
flow close to the surface. Similar riblets can reduce friction drag in selected conditions, although
their benefit depends on size, orientation and cleanliness.

Nature also inspires structures. Bone combines a dense outer layer with a porous interior, and achieves a
high stiffness-to-mass ratio. Additive manufacturing can produce lattice structures that use the same
principle. Honeycomb panels and branching networks similarly offer efficient ways to distribute
loads or fluids.

A biomimetic design is not automatically sustainable. Manufacturing a complex microstructure may
consume large amounts of energy or use materials that are difficult to recycle. Engineers must compare
the complete life cycle, durability and maintenance requirements with a conventional solution.
Biomimetics is most valuable when the biological analogy leads to a measurable improvement rather
than merely providing an attractive product story.
\end{ecarticle}

\begin{tcolorbox}[title=\sffamily\bfseries Did You Know?,colback=yellow!10,colframe=ECOrange]
The hook-and-loop fastener was inspired by burrs that attached themselves to animal fur and clothing.
The invention reproduced the attachment principle, not the complete plant structure.
\end{tcolorbox}

\begin{engineeringfocus}
\textbf{From biological observation to engineering validation}
\begin{center}
\resizebox{\linewidth}{!}{%
\begin{tikzpicture}[font=\scriptsize,>=Latex,node distance=3.3mm]
\node[draw=ECPurple,fill=ECPurple!10,rounded corners,align=center] (p) {define the\\technical function};
\node[draw=ECGreen,fill=ECGreen!10,rounded corners,align=center,right=of p] (b) {observe a\\biological strategy};
\node[draw=ECBlue,fill=ECLightBlue,rounded corners,align=center,right=of b] (a) {abstract the\\working principle};
\node[draw=ECOrange,fill=ECOrange!10,rounded corners,align=center,below=of a] (m) {model and\\prototype};
\node[draw=ECRed,fill=ECRed!8,rounded corners,align=center,below=of b] (v) {test performance\\and life cycle};
\draw[->,thick] (p)--(b); \draw[->,thick] (b)--(a); \draw[->,thick] (a)--(m);
\draw[->,thick] (m)--(v); \draw[->,thick,dashed] (v)--(p);
\end{tikzpicture}%
}
\end{center}
Iteration is essential because biological materials and manufacturing constraints are very different.
\end{engineeringfocus}

\begin{keyfigures}
\begin{tabularx}{\linewidth}{@{}Xr@{}}
Primary objective & functional transfer\\
Common scale & microstructure\\
Typical benefit & low mass or drag\\
Necessary step & abstraction\\
Final proof & measured validation
\end{tabularx}
\end{keyfigures}

\begin{tcolorbox}[title=\sffamily\bfseries Illustrative drag comparison,colback=white,colframe=ECPurple]
\begin{center}
\begin{tikzpicture}
\begin{axis}[xbar,width=0.97\linewidth,height=49mm,xlabel={Relative drag index},xmin=0,xmax=110,
symbolic y coords={Optimised riblets,Shark-inspired coating,Riblet texture,Smooth reference},
ytick=data,grid=major,ticklabel style={font=\scriptsize},nodes near coords]
\addplot table[x=drag,y=surface,col sep=comma]{data/EC19/drag.csv};
\end{axis}
\end{tikzpicture}
\end{center}
\footnotesize Illustrative values showing that geometry must be optimised for the operating flow.
\end{tcolorbox}

\begin{vocabulary}
\begin{tabularx}{\linewidth}{@{}lX@{}}
biomimetics & biomimétisme\\
living system & système vivant\\
working principle & principe de fonctionnement\\
hierarchical structure & structure hiérarchique\\
riblet & micro-rainure\\
adhesion & adhérence\\
porous & poreux\\
lattice structure & structure en treillis\\
stiffness-to-mass ratio & rapport rigidité/masse\\
life cycle & cycle de vie
\end{tabularx}
\end{vocabulary}


\begin{applicationexercise}
A smooth reference surface has a drag index of 100. A biomimetic riblet surface has an index of 82,
but its manufacture increases the component mass from 12.0 kg to 12.6 kg.
\begin{enumerate}
\item Calculate the percentage reduction in drag index.
\item Calculate the percentage increase in mass.
\item A design team assigns equal importance to drag reduction and mass increase. Compare the two
percentages and decide whether the riblet solution appears favourable according to this simplified
criterion.
\item State two tests that would be required before approving the design for real service.
\end{enumerate}
\end{applicationexercise}

\begin{questions}
\item Why is direct imitation often insufficient?
\item Explain how gecko feet inspire dry adhesives.
\item What engineering benefit can be obtained from bone-like structures?
\item Why must operating scale be considered for riblets?
\item How can a biomimetic product have a poor environmental performance?
\end{questions}

\begin{engineeringchallenge}
Choose a biological strategy that could improve a small wind turbine, robot gripper or protective
helmet. Identify the function, abstract the principle, propose a manufacturable concept and define
two experiments that compare it with a conventional design.
\end{engineeringchallenge}

\begin{speaking}
Is nature always the best engineer? Discuss optimisation, evolution, materials, manufacturing and the
need for measurable evidence.
\end{speaking}
\end{multicols}

\subsection*{Sources}
\ECsource{AskNature -- Biomimicry examples}{https://asknature.org/}
\ECsource{Biomimicry Institute}{https://biomimicry.org/}
